\documentclass[a4paper, 12pt, landscape]{scrartcl}

\usepackage[a4paper, landscape, margin=1.0cm]{geometry}
\usepackage[ngerman]{babel}
\usepackage{libertinus}
\usepackage{xcolor}
\usepackage[most]{tcolorbox}

\pagestyle{empty}

\definecolor{primaryBlue}{RGB}{20, 50, 100}
\definecolor{cardBg}{RGB}{252, 253, 255}

\newcommand{\kernaussageStreifen}[4]{%
    \noindent\begin{tcolorbox}[
        enhanced,
        colback=cardBg,
        colframe=primaryBlue,
        boxrule=3pt,
        arc=4mm,
        valign=center,
        halign=center,
        height=17.5cm,
        top=0.6cm, bottom=0.6cm, left=1.0cm, right=1.0cm,
        title={\LARGE\bfseries #1 \enskip$\cdot$\enskip \normalsize #2},
        fonttitle=\bfseries\color{white},
        attach boxed title to top center={yshift=-4.5mm},
        boxed title style={colback=primaryBlue, arc=2mm, boxrule=0pt, left=10mm, right=10mm, top=3mm, bottom=3mm}
    ]
        \vspace{0.4cm}
        {\fontsize{38pt}{48pt}\selectfont\sffamily\bfseries\color{primaryBlue} \glqq #3\grqq}
        
        \vfill
        {\large\bfseries\color{gray} #4}
    \end{tcolorbox}
    \newpage
}

\begin{document}

% --- SATZ A ---
\kernaussageStreifen{SATZ A}{Kernaussage}{Das Tier produziert nur unter der Herrschaft des physischen Bedürfnisses, während der Mensch frei produziert und erst wahrhaft produziert in der Freiheit [\dots] der Mensch formiert daher auch nach den Gesetzen der Schönheit.}{Karl Marx: Ökonomisch-philosophische Manuskripte (1844)}

% --- SATZ B ---
\kernaussageStreifen{SATZ B}{Kernaussage}{Was das Produkt seiner Arbeit ist, ist er nicht. Je mehr Gegenstände der Arbeiter produziert, um so mehr gerät er unter die Herrschaft seines Produkts, des Kapitals.}{Karl Marx: Ökonomisch-philosophische Manuskripte (1844)}

% --- SATZ C ---
\kernaussageStreifen{SATZ C}{Kernaussage}{Seine Arbeit ist daher nicht freiwillig, sondern gezwungen, Zwangsarbeit. Sie ist nicht die Befriedigung eines Bedürfnisses, sondern nur ein Mittel, um Bedürfnisse außer ihr zu befriedigen.}{Karl Marx: Ökonomisch-philosophische Manuskripte (1844)}

% --- SATZ D ---
\kernaussageStreifen{SATZ D}{Kernaussage}{Es kommt daher zum Resultat, dass der Mensch sich nur mehr in seinen tierischen Funktionen (Essen, Trinken, Zeugen) als freitätig fühlt und in seinen menschlichen Funktionen nur mehr als Tier.}{Karl Marx: Ökonomisch-philosophische Manuskripte (1844)}

\end{document}
